\documentclass[10pt,a4paper]{amsart}
\usepackage[margin=19mm]{geometry}
\usepackage{amsmath,amssymb,mathtools}
\usepackage[scr=rsfs,cal=euler]{mathalfa}
\usepackage{xfrac}
\usepackage[unicode,hidelinks,bookmarks=false]{hyperref}
\newcommand{\ie}{i.e.\ }
\newcommand{\cf}{cf.\ }
\newcommand{\resp}{respectively\ }
\newcommand{\Subsection}{\subsection}
\providecommand{\nobreakdash}{\nobreak\hbox{-}}
\setlength{\emergencystretch}{2em}
\multlinegap0pt
\usepackage{amssymb}
\usepackage{tikz}
\newtheorem{theorem}{Theorem}
\newtheorem{letlist}{Letlist}
\renewcommand\a{\alpha}
\newcommand\g{\gamma}
\newcommand\oo{\infty}
\newcommand\s{\sigma}
\title{Self-avoiding walks on cubic graphs and local transformations}
\author{Benjamin Grant, Zhongyang Li}
\date{Source: arXiv:2601.12571v2 (CC BY 4.0)}
\begin{document}
\maketitle

\noindent\emph{Source and licence.} Self-avoiding walks on cubic graphs and local transformations. Version arXiv:2601.12571v2 (CC BY 4.0), distributed by arXiv under the Creative Commons Attribution 4.0 International licence, http://creativecommons.org/licenses/by/4.0/, as recorded in the arXiv licence metadata for this exact version and shown on the abstract page https://arxiv.org/abs/2601.12571v2. Full licence notice: \texttt{licenses/arxiv-2601.12571-CC-BY-4.0.txt}.\par\smallskip

\noindent\textit{Editorial scope.} Counted unit: the article's theorem thm:main2' (source0.tex lines 508--519). The article's complete original proof of it is delivered with it (source0.tex lines 671--1012, one \texttt{proof} environment of 342 lines, which the article places away from the statement). No mathematics is written, repaired or re-derived: the statement and the proof are the article's own lines, in the article's own notation, and no retained line is substituted or re-flowed.  Retained uncounted as the source-local context the proof needs: lines 458--460; lines 463--466; lines 486--491.  Nothing was omitted from any retained span, so the package carries no omission record.  No retained line contains a citation, so the wrapper carries no bibliography and no bibliography entry.  No retained line contains a figure environment, an image-inclusion command or drawing code, so no image is delivered and the package declares no asset dependency.  Editorial formatting only, in addition: an AMS \texttt{amsart} preamble that restates the article's own declarations and macros used by the retained lines, namely the environments \texttt{theorem}, \texttt{letlist} on separate counters, together with the standard packages those lines need.  The retained environments print in this document as Theorem 1, Letlist 1 in order of appearance, whereas the article prints the same items with the numbering of its own full document.  The complete native source of the article as distributed in the arXiv e-print for this exact version is bundled unchanged as \texttt{sources/192961/source0.tex}.
\begin{equation}\label{def:slow}
C_1 L(n)n^{\gamma-1}\mu^n\leq \sigma_{n}\leq C_2 L(n)n^{\gamma-1}\mu^n, \qquad n \ge 1.
\end{equation}

\begin{equation}\label{def0}
V(z)=\sum_{n=1}^{\infty}\frac{1}{n^{2z+1}} \langle \|\pi\|^2 \rangle_n,
\qquad z \in [-\oo,\oo],
\end{equation}

\begin{equation}\label{dd1}
Z_\a(\mu^{-1})d_G(W,\a W)^{w} 
\begin{cases} \to 0 &\text{if } w<\eta,\\
\to\infty &\text{if } w>\eta. 
\end{cases}
\end{equation}

\begin{theorem}\label{thm:main2'}
Let $G_0$ be an infinite, quasi-transitive, connected, cubic graph.
Assume that $|\g(G_0)|<\oo$ and that $\eta(G_1)$ exists.
\begin{letlist}
\item The exponents $\g$, $\eta$ of $G_0$ and $G_1$ are equal.
\item  Let $\sigma_{n,k}$ be
the number of $n$-step SAWs on $G_k$ from mid-edges in $X_k$.
Assume the $\s_{n,k}$ satisfy \eqref{def:slow} for constants
$C_{i,k}$ and a common slowly varying function $L$. Then
the exponents $\nu$ of $G_0$ and $G_1$ are equal.
\end{letlist}
\end{theorem}

\begin{proof}[Proof of Theorem~\ref{thm:main2'}]
Throughout, let $G_0:=G$ and $G_1:=\phi(G)$, and let $W_k,X_k,\Sigma_k$ be as in
Section~3.  Let $\Sigma_1^\ast$ be the set of SAWs on $G_1$ that start at mid-edges of $X_0$
and end at mid-edges of $E_0$ (viewed as mid-edges of $E_1$).  For $n\ge 0$, write
$\Sigma_{n,k}$ (resp.\ $\Sigma_{n,1}^\ast$) for the subset of $\Sigma_k$ (resp.\ $\Sigma_1^\ast$)
containing $n$-step SAWs, and let $\sigma_{n,k}:=|\Sigma_{n,k}|$ and $\sigma_{n,1}^\ast:=|\Sigma_{n,1}^\ast|$.

Let $N:=|V(\phi(v))|$ be the number of vertices of a gadget (independent of $v$ by
Definition~3.1(2)).  Let $\Delta:=\mathrm{diam}(\phi(v))$ and $D:=\max_{u\in V(\phi(v))}\deg_{\phi(v)}(u)$
(both independent of $v$).  Note that any SAW segment inside a gadget connecting two
ports has length between $2$ and $N$ (in the mid-edge convention of Section~2).

\medskip\noindent
\textbf{(a) Equality of $\gamma$ and $\eta$.}
For $k=0,1$ define
\[
Y_k(x,y)\;=\;\sum_{\pi\in\Sigma_k}\frac{x^{|\pi|}}{|\pi|^{\,y}},
\qquad
Y_1^\ast(x,y)\;=\;\sum_{\pi\in\Sigma_1^\ast}\frac{x^{|\pi|}}{|\pi|^{\,y}},
\qquad x>0,\ y\in\mathbb{R},
\]
with the convention that the denominator is $1$ when $|\pi|=0$.

\smallskip\noindent
\emph{Step 1: $Y_1^\ast(x,y)<\infty \iff Y_1(x,y)<\infty$ for fixed $(x,y)$.}
Since $\Sigma_1^\ast\subseteq \Sigma_1$, we have $Y_1^\ast(x,y)\le Y_1(x,y)$.

Conversely, consider $\pi\in\Sigma_1\setminus\Sigma_1^\ast$.  Either:
(i) $\pi$ is entirely contained in a single gadget (there are only finitely many such walks,
of bounded length $\le N$), or
(ii) $\pi$ visited at least one mid-edge in $E_0$ and at least one endpoint of $\pi$ is a mid-edge not in $E_0$.
In case (ii), move each endpoint that is not in $E_0$ along $\pi$ to the first mid-edge of $E_0$
encountered when traversing $\pi$ from that endpoint.  This removes at most $N$ visited vertices
per endpoint, hence produces a (strictly) shorter SAW $\pi^\dagger\in \Sigma_1^\ast$ with
$|\pi|-2N\le |\pi^\dagger|\le |\pi|$.

Moreover, for a fixed $\pi^\dagger$, the number of possible $\pi$ that reduce to it by this procedure
is bounded by a constant depending only on $(N,D)$ (each endpoint extension has at most $\sum_{j=0}^{N}(Dx)^j$
choices at weight $x$, and there are two endpoints).  Since $|\pi|$ and $|\pi^\dagger|$ differ by at most $2N$,
we also have for fixed $y\in\mathbb{R}$ a constant $C_y<\infty$ such that
$|\pi|^{-y}\le C_y\,|\pi^\dagger|^{-y}$.
It follows that for each fixed $(x,y)$,
\[
Y_1(x,y)\;\le\;C(x,y)\,Y_1^\ast(x,y)\;+\;C'(x,y),
\]
where $C'(x,y)<\infty$ accounts for the finitely many SAWs contained in a single gadget.
Therefore,
\begin{equation}\label{eq:Ystar-equiv}
Y_1^\ast(x,y)<\infty \iff Y_1(x,y)<\infty .
\end{equation}

\smallskip\noindent
\emph{Step 2: compare $Y_1^\ast(x,y)$ and $Y_0(g(x),y)$.}
Fix an $n$-step SAW $\omega\in\Sigma_{n,0}$.  Contracting each gadget of $G_1$ to a vertex
maps every $\pi\in\Sigma_1^\ast$ to a unique $\omega\in\Sigma_0$; conversely, expanding each visit
to a vertex of $\omega$ by a SAW segment inside the corresponding gadget yields SAWs in $\Sigma_1^\ast$.
Let $a_m:=\sigma^{\phi(v)}_{m}(e_1,e_2)$ be the number of $m$-step gadget SAWs (in the mid-edge convention)
connecting two ports, so that
\[
g(x)=\sum_{m\ge 0} a_m x^{m},
\qquad a_0=a_1=0,\qquad a_m=0 \text{ for } m>N.
\]
Let $b_{n,m}$ be the coefficient of $x^m$ in $g(x)^n$, i.e.\ $g(x)^n=\sum_m b_{n,m}x^m$.
Then $b_{n,m}\neq 0$ only when $2n\le m\le Nn$, and the total contribution to $Y_1^\ast(x,y)$
from all $\pi\in\Sigma_1^\ast$ that contract to this fixed $\omega$ equals
\[
T_n(x,y)\;:=\;\sum_{m=2n}^{Nn} b_{n,m}\,\frac{x^{m}}{m^{\,y}}.
\]
Since $2n\le m\le Nn$, we have
\[
(2n)^{-y}\wedge (Nn)^{-y}\;\le\; m^{-y}\;\le\;(2n)^{-y}\vee (Nn)^{-y}.
\]
Hence, with
\[
c_y := 2^{-y}\wedge N^{-y},
\qquad
c_y' := 2^{-y}\vee N^{-y},
\]
we obtain the bounds
\[
c_y\,\frac{g(x)^n}{n^{\,y}}
\;\le\;
T_n(x,y)
\;\le\;
c_y'\,\frac{g(x)^n}{n^{\,y}}.
\]
Summing $T_n(x,y)$ over all $\omega\in\Sigma_{n,0}$ and then over $n\ge 1$ gives
\begin{equation}\label{eq:Ystar-Y0-compare}
c_y\,\widetilde Y_0(g(x),y)\;\le\;\widetilde Y_1^\ast(x,y)\;\le\;c_y'\,\widetilde Y_0(g(x),y),
\end{equation}
where $\widetilde Y$ denotes the series with the (harmless) $n=0$ term removed.
In particular,
\begin{equation}\label{eq:Ystar-finite-iff}
Y_1^\ast(x,y)<\infty \iff Y_0(g(x),y)<\infty .
\end{equation}

By Theorem~3.3, $\mu_0^{-1}=g(\mu_1^{-1})$.  Combining \eqref{eq:Ystar-equiv} and \eqref{eq:Ystar-finite-iff}
at $x=\mu_1^{-1}$ yields
\[
Y_1(\mu_1^{-1},y)<\infty \iff Y_0(\mu_0^{-1},y)<\infty,
\]
and therefore $\gamma(G_1)=\gamma(G_0)$.

\smallskip\noindent
\emph{Step 3: equality of $\eta$.}
For $\alpha\in A(G_0)$, and $k\in\{0,1\}$ define 
\begin{align*}
Z_{\alpha,k}(x)=\sum_{\pi\in \Sigma_k(\alpha)}x^{|\pi|}
\end{align*}
where $\Sigma_k(\alpha
)$ is the subset of $\Sigma_k$ containing SAWs ending at mid-edges in $\alpha X_k$.
Define $Z_{\alpha,1}^\ast(x)$ analogously but with SAWs on $G_1$ from mid-edges of $X_0$ to mid-edges of
$\alpha X_0$ (both viewed in $G_1$).  As in the proof of Theorem~3.3, endpoint-adjustment inside gadgets gives
a two-sided bound
\begin{equation}\label{eq:Zalpha-compare}
Z_{\alpha,1}^\ast(x)\;\le\; Z_{\alpha,1}(x)\;\le\;C(x)\,Z_{\alpha,1}^\ast(x),
\end{equation}
for a finite constant $C(x)$ depending only on the gadget and $x$.
Moreover, shrinking gadgets yields the exact identity
\begin{equation}\label{eq:Zalpha-subst}
Z_{\alpha,1}^\ast(x)\;=\;Z_{\alpha,0}(g(x)).
\end{equation}
Evaluating at $x=\mu_1^{-1}$ and using $\mu_0^{-1}=g(\mu_1^{-1})$ gives
$Z_{\alpha,1}^\ast(\mu_1^{-1})=Z_{\alpha,0}(\mu_0^{-1})$.

Finally, graph distances in $G_0$ and $G_1$ are comparable up to multiplicative constants:
there exists $\kappa\in(0,\infty)$ (depending only on the gadget) such that for all $\alpha$,
\[
d_{G_0}(W_0,\alpha W_0)\;\le\; d_{G_1}(W_1,\alpha W_1)\;\le\;\kappa\, d_{G_0}(W_0,\alpha W_0).
\]
Using \eqref{eq:Zalpha-compare}--\eqref{eq:Zalpha-subst} and the above distance comparison in the definition
\eqref{dd1} of $\eta$, we obtain $\eta(G_1)=\eta(G_0)$ (under the assumed existence of $\eta(G_1)$).
This completes the proof of part~(a).

\medskip\noindent
\textbf{(b) Equality of $\nu$ under \eqref{def:slow}.}
For $k=0,1$, let
\[
V_k(z)\;=\;\sum_{n=1}^\infty \frac{1}{n^{2z+1}}\,\Big\langle \|\pi\|_k^{\,2}\Big\rangle_{n,k},
\]
as in \eqref{def0}, and define
\begin{align}
V_1^\ast(z)\;:=\;\sum_{n=1}^\infty \frac{1}{n^{2z+1}}\,
\frac{1}{\sigma_{n,1}}\sum_{\pi\in\Sigma_{n,1}^\ast}\|\pi\|_1^{\,2}.\label{dv1s}
\end{align}
(Equivalently, $V_1^\ast(z)=\sum_{n\ge 1}n^{-(2z+1)}(\sigma_{n,1}^\ast/\sigma_{n,1})\langle\|\pi\|_1^2\rangle_{n,1}^\ast$.)

\smallskip\noindent
\emph{Step 4: $V_1(z)<\infty \iff V_1^\ast(z)<\infty$ for each finite $z$.}
Clearly $V_1^\ast(z)\le V_1(z)$.
For the reverse implication,  



Recall that
\[
V_1(z)=\sum_{n\ge1}\frac{1}{n^{2z+1}}\Big\langle \|\pi\|_1^2\Big\rangle_{n,1}
=\sum_{n\ge1}\frac{1}{n^{2z+1}\sigma_{n,1}}\sum_{\pi\in\Sigma_{n,1}}\|\pi\|_1^2,
\]
and, by \eqref{dv1s},
\[
V_1^\ast(z)=\sum_{n\ge1}\frac{1}{n^{2z+1}}\frac{\sigma_{n,1}^\ast}{\sigma_{n,1}}
\Big\langle \|\pi\|_1^2\Big\rangle_{n,1}^\ast
=\sum_{n\ge1}\frac{1}{n^{2z+1}\sigma_{n,1}}\sum_{\pi\in\Sigma_{n,1}^\ast}\|\pi\|_1^2.
\]
Hence
\begin{equation}\label{eq:V1-decomp}
V_1(z)=V_1^\ast(z)+R(z),
\qquad
R(z):=\sum_{n\ge1}\frac{1}{n^{2z+1}\sigma_{n,1}}
\sum_{\pi\in\Sigma_{n,1}\setminus\Sigma_{n,1}^\ast}\|\pi\|_1^2.
\end{equation}

\smallskip\noindent
\textbf{(i) Walks contained in a single gadget give a finite additive constant.}
Let $\mathcal{G}$ be the set of SAWs in $\Sigma_{1}$ that are entirely contained in a single
gadget. Since the walk starts at a mid-edge in $X_1$, only finitely many gadgets intersect
the fundamental domain $W_1$, and each gadget is finite. Therefore $\mathcal{G}$ is finite,
and its total contribution
\[
C_z' := \sum_{\pi\in\mathcal{G}}
\frac{\|\pi\|_1^2}{|\pi|^{2z+1}\sigma_{|\pi|,1}}
\]
is finite for each fixed finite $z$.

\smallskip\noindent
\textbf{(ii) Trimming map to $\Sigma_{1}^\ast$ with bounded multiplicity.}
 For $\pi\in\Sigma_{n,1}\setminus\Sigma_{n,1}^\ast$ that is not
contained in a single gadget, at least one endpoint of $\pi$ is not a mid-edge of $E_0$.
Move each endpoint along $\pi$ to the first mid-edge of $E_0$ encountered, deleting the
initial segment at that endpoint. This produces a shorter SAW $\tau(\pi)\in\Sigma_{m,1}^\ast$
with
\[
m=|\tau(\pi)|\in\{n-j:\ 0\le j\le 2N\}.
\]
Moreover, %for each fixed $j\in\{0,\dots,2N\}$, the map
%\[
%\tau_j:\ \{\pi\in\Sigma_{n,1}\setminus\Sigma_{n,1}^\ast:\ |\pi|-|\tau(\pi)|=j\}\to\Sigma_{n-j,1}^\ast
%\]
%has uniformly bounded multiplicity: there exists a constant
\[
K:=\Bigl(1+\sum_{i=1}^{N}D^{i}\Bigr)^2<\infty
\]
such that every $\pi^\ast\in\Sigma_{m,1}^\ast$ has at most $K$ preimages (each endpoint can be
extended by at most $N$ steps, each step has at most $D$ choices; there are two endpoints) under $\tau$.

Also, trimming changes each endpoint by at most $N$ vertices within a gadget, hence changes
the endpoint distance by at most an additive constant. In particular,
\[
\|\pi\|_1 \le \|\tau(\pi)\|_1 + 2M.
\]
Since $|\tau(\pi)|\ge 1$ implies $\|\tau(\pi)\|_1\ge 1$, we have
\begin{equation}\label{eq:dist-bound}
\|\pi\|_1^2 \le (1+2M)^2\,\|\tau(\pi)\|_1^2 .
\end{equation}

\smallskip\noindent
\textbf{(iii) Comparing the normalizing factors using \eqref{def:slow}.}
Fix $j\in\{0,\dots,2N\}$ and finite $z$. For $m\ge1$,
\begin{equation}\label{eq:weight-compare}
\frac{1}{(m+j)^{2z+1}\sigma_{m+j,1}}
=
\Bigl(\frac{m}{m+j}\Bigr)^{2z+1}\frac{\sigma_{m,1}}{\sigma_{m+j,1}}
\cdot \frac{1}{m^{2z+1}\sigma_{m,1}}.
\end{equation}
By \eqref{def:slow} (i.e.\ (3.7)) for $k=1$,
\[
\sigma_{n,1}\asymp L(n)\,n^{\gamma-1}\mu_1^{\,n},
\]
so for fixed $j$,
\[
\frac{\sigma_{m,1}}{\sigma_{m+j,1}}
\le
\frac{C_{2,1}}{C_{1,1}}\,
\frac{L(m)}{L(m+j)}\,
\Bigl(\frac{m}{m+j}\Bigr)^{\gamma-1}\,
\mu_1^{-j}.
\]
Since $L$ is slowly varying, $L(m)/L(m+j)$ is bounded in $m$ for each fixed $j$,
and $\frac{m}{m+j}\in[1/(1+j),1]$ so the power term is bounded as well. Therefore the product
\[
A_{z,j}:=\sup_{m\ge1}\Bigl(\frac{m}{m+j}\Bigr)^{2z+1}\frac{\sigma_{m,1}}{\sigma_{m+j,1}}<\infty,
\]
and \eqref{eq:weight-compare} yields
\begin{equation}\label{eq:factor-bound}
\frac{1}{(m+j)^{2z+1}\sigma_{m+j,1}}
\le A_{z,j}\cdot \frac{1}{m^{2z+1}\sigma_{m,1}}.
\end{equation}

\smallskip\noindent
\textbf{(iv) Bounding $R(z)$ by $V_1^\ast(z)$.}
Using \eqref{eq:dist-bound}, bounded multiplicity $K$, and \eqref{eq:factor-bound}, we obtain
\begin{align*}
R(z)
&\le C_z' + \sum_{j=0}^{2N}\sum_{m\ge1}
\frac{1}{(m+j)^{2z+1}\sigma_{m+j,1}}
\sum_{\substack{\pi\in\Sigma_{m+j,1}\setminus\Sigma_{m+j,1}^\ast\\ |\pi|-|\tau(\pi)|=j}}
\|\pi\|_1^2 \\
&\le C_z' + \sum_{j=0}^{2N}\sum_{m\ge1}
\frac{1}{(m+j)^{2z+1}\sigma_{m+j,1}}
\sum_{\pi^\ast\in\Sigma_{m,1}^\ast} K\,(1+2N)^2 \|\pi^\ast\|_1^2 \\
&\le C_z' + \sum_{j=0}^{2N} K\,(1+2N)^2 A_{z,j}
\sum_{m\ge1}\frac{1}{m^{2z+1}\sigma_{m,1}}\sum_{\pi^\ast\in\Sigma_{m,1}^\ast}\|\pi^\ast\|_1^2 \\
&= C_z' + C_z\, V_1^\ast(z),
\end{align*}
where $C_z:=\sum_{j=0}^{2N} K(1+2N)^2A_{z,j}<\infty$.
Combining this with \eqref{eq:V1-decomp} gives the explicit bound
\begin{equation}\label{eq:V1-upper-explicit}
V_1(z)\le (1+C_z)\,V_1^\ast(z)+C_z'.
\end{equation}
Since also $V_1^\ast(z)\le V_1(z)$, we conclude that for each fixed finite $z$,
\[
V_1(z)<\infty \iff V_1^\ast(z)<\infty,
\]
which is \eqref{eq:V1-V1star}.





Therefore \begin{equation}\label{eq:V1-V1star} V_1(z)\;<\infty \iff V_1^\ast(z)\;<\infty .
\end{equation}


\smallskip\noindent
\emph{Step 5: $V_1^\ast(z)<\infty \iff V_0(z)<\infty$ for each finite $z$.}
Fix $n\ge 1$ and $\omega\in\Sigma_{n,0}$.  Let $\mathcal{E}(\omega)$ be the family of SAWs
$\pi\in\Sigma_1^\ast$ that contract to $\omega$.  As above, for each $m$ the number of $\pi\in\mathcal{E}(\omega)$
with $|\pi|=m$ equals $b_{n,m}$ (the coefficient of $x^m$ in $g(x)^n$), and $b_{n,m}=0$ unless $2n\le m\le Nn$.

Since endpoints of $\pi$ and $\omega$ correspond to the same two original mid-edges, the graph distances satisfy
$\|\omega\|_0 \le \|\pi\|_1 \le \kappa \|\omega\|_0$ for some $\kappa<\infty$ depending only on the gadget.
Hence, for fixed $\omega$,
\[
\sum_{\pi\in\mathcal{E}(\omega)} \frac{\|\pi\|_1^{\,2}}{|\pi|^{2z+1}\,\sigma_{|\pi|,1}}
\asymp
\|\omega\|_0^{\,2}\sum_{m=2n}^{Nn} b_{n,m}\,\frac{1}{m^{2z+1}\,\sigma_{m,1}},
\]
where $\asymp$ hides multiplicative constants depending only on $(z,\kappa)$.

Now apply \eqref{def:slow} (for $k=1$): for all $m\ge 1$,
\[
\frac{1}{\sigma_{m,1}}
\asymp
\frac{\mu_1^{-m}}{L(m)\,m^{\gamma-1}}.
\]
Therefore,
\[
\sum_{m=2n}^{Nn} b_{n,m}\,\frac{1}{m^{2z+1}\,\sigma_{m,1}}
\asymp
\sum_{m=2n}^{Nn} b_{n,m}\,\frac{\mu_1^{-m}}{L(m)\,m^{2z+\gamma}}.
\]
Since $L$ is slowly varying and $m\in[2n,Nn]$, the values $L(m)$ are comparable to $L(n)$ up to a constant
(depending on $N$).  Also $m^{-(2z+\gamma)}$ is comparable to $n^{-(2z+\gamma)}$ up to a constant depending
on $(z,N)$.  Hence
\[
\sum_{m=2n}^{Nn} b_{n,m}\,\frac{\mu_1^{-m}}{L(m)\,m^{2z+\gamma}}
\asymp
\frac{1}{L(n)\,n^{2z+\gamma}}\sum_{m=2n}^{Nn} b_{n,m}\,\mu_1^{-m}.
\]
But $\sum_m b_{n,m}\mu_1^{-m}=g(\mu_1^{-1})^n=\mu_0^{-n}$ by Theorem~3.3.
Thus,
\[
\sum_{\pi\in\mathcal{E}(\omega)} \frac{\|\pi\|_1^{\,2}}{|\pi|^{2z+1}\,\sigma_{|\pi|,1}}
\asymp
\|\omega\|_0^{\,2}\,\frac{\mu_0^{-n}}{L(n)\,n^{2z+\gamma}}.
\]
Summing over $\omega\in\Sigma_{n,0}$ and then over $n\ge 1$ yields
\[
V_1^\ast(z)\asymp \sum_{n\ge 1}\frac{1}{\sigma_{n,0}}\frac{1}{n^{2z+1}}\sum_{\omega\in\Sigma_{n,0}}\|\omega\|_0^{\,2}
\;=\;V_0(z),
\]
where we used again \eqref{def:slow} for $k=0$ to identify
$\sigma_{n,0}^{-1}n^{-(2z+1)}\asymp \mu_0^{-n}(L(n)n^{2z+\gamma})^{-1}$.
Consequently, for finite $z$,
\begin{equation}\label{eq:V0-V1star}
V_1^\ast(z)<\infty \iff V_0(z)<\infty .
\end{equation}

Combining \eqref{eq:V1-V1star} and \eqref{eq:V0-V1star}, we conclude that
$V_1(z)<\infty \iff V_0(z)<\infty$ for all finite $z$, hence $\nu(G_1)=\nu(G_0)$.
\end{proof}

\end{document}
